\documentclass[a4paper,oneside,11pt]{book}
\usepackage{graphicx}
\usepackage{titling}
\usepackage{geometry}
\usepackage{listings}
\usepackage{xcolor}
\usepackage{float}
\usepackage{indentfirst}
\usepackage{longtable}
\usepackage[T1]{fontenc}
\usepackage{lmodern}
\usepackage[hyphens]{url}
\usepackage{hyperref}
\def\UrlBreaks{\do\/\do-}

\geometry{left=3cm, right=2.5cm, top=3cm, bottom=3cm}
\emergencystretch=5em

\definecolor{codegreen}{rgb}{0,0.6,0}
\definecolor{codegray}{rgb}{0.5,0.5,0.5}
\definecolor{codepurple}{rgb}{0.58,0,0.82}
\definecolor{backcolour}{rgb}{0.97,0.97,0.95}
\definecolor{codeblue}{rgb}{0.16,0.32,0.75}

\lstdefinelanguage{yaml}{
    keywords={true,false,null,yes,no,on,needs,if,runs-on,uses,with,run,steps,jobs,name,env,environment},
    keywordstyle=\color{codeblue}\bfseries,
    basicstyle=\ttfamily\footnotesize,
    sensitive=false,
    comment=[l]{\#},
    commentstyle=\color{codegray}\itshape,
    stringstyle=\color{codepurple},
    morestring=[b]',
    morestring=[b]"
}

\lstdefinestyle{codestyle}{
    backgroundcolor=\color{backcolour},
    basicstyle=\ttfamily\footnotesize,
    commentstyle=\color{codegreen},
    keywordstyle=\color{codeblue}\bfseries,
    stringstyle=\color{codepurple},
    numberstyle=\tiny\color{codegray},
    numbers=left,
    numbersep=6pt,
    showspaces=false,
    showstringspaces=false,
    breaklines=true,
    tabsize=2,
    frame=single,
    rulecolor=\color{gray},
    captionpos=b,
}
\lstset{style=codestyle}

\title{Laporan Praktikum \\ Konstruksi dan Evolusi Perangkat Lunak\\
Pertemuan 3 \\ CI/CD Pipeline Empat Tahap dengan Laravel}
\author{Hisyam Al Rayyan / 24/541582/SV/24926}
\date{\today}

\begin{document}

\begin{titlingpage}
\begin{center}
\vspace{4cm}
\begin{huge}
\textbf{\thetitle} \\
\end{huge}
\vspace{2cm}
\includegraphics[height=8cm]{images/logo_ugm.png}\\
\vspace{4cm}
{\Large
\begin{tabular}{rl}
Nama            & : Hisyam Al Rayyan \\
NIM             & : 24/541582/SV/24926 \\
Kelas           & : A2 \\
Dosen Pengampu  & : Galih Malela Damaraji \\
\end{tabular}
}
\vspace{2cm}
\end{center}
\end{titlingpage}

\tableofcontents

%=== BAB 1 ===
\chapter{Pendahuluan}

\section{Tujuan Praktikum}
\begin{enumerate}
    \item Memahami konsep pipeline CI/CD empat tahap: \textit{build}, \textit{test}, \textit{staging}, dan \textit{production}.
    \item Menerapkan dependensi antar-job menggunakan kata kunci \texttt{needs:} pada GitHub Actions.
    \item Membuat fitur CRUD sederhana dengan Laravel sebagai dasar pengujian pipeline.
    \item Mengamankan job \textit{production} agar hanya berjalan dari branch \texttt{main} dengan mekanisme \textit{environment} dan \textit{required reviewer}.
    \item Menyusun skrip \texttt{deploy.sh} dengan tujuh langkah deploy.
    \item Membuktikan bahwa pipeline berhenti ketika pengujian gagal dan \textit{production} dilewati pada branch selain \texttt{main}.
\end{enumerate}

\section{Alat dan Bahan}
\begin{enumerate}
    \item Laptop dengan sistem operasi Windows.
    \item PHP 8.4 dan Composer.
    \item Git dan GitHub CLI (\texttt{gh}).
    \item Repository GitHub \texttt{evolusi-pl-24-541582-SV-24926} pada organization KEPL2026.
    \item Framework Laravel 13.
    \item GitHub Actions sebagai platform CI/CD.
\end{enumerate}

%=== BAB 2 ===
\chapter{Hasil dan Pembahasan}

\section{Pembuatan Fitur CRUD Tugas}

Sebagai fondasi pengujian pipeline, dibuat fitur CRUD sederhana untuk entitas \texttt{Tugas}. Tabel \texttt{tugas} memiliki kolom \texttt{id}, \texttt{judul}, \texttt{deskripsi}, \texttt{status} (enum: \textit{pending}/\textit{selesai}), dan \texttt{timestamps}. Migration, model, dan controller dibuat menggunakan perintah Artisan Laravel.

\begin{lstlisting}[language=bash, caption={Perintah pembuatan CRUD Tugas}]
php artisan make:model Tugas -mcr
\end{lstlisting}

Model \texttt{Tugas} didefinisikan dengan properti \texttt{fillable} agar mass assignment dapat dilakukan. Controller \texttt{TugasController} mengimplementasikan operasi CRUD lengkap. Endpoint JSON disediakan melalui \texttt{routes/api.php} pada path \texttt{GET /api/tugas}.

% Screenshot 1: Struktur CRUD
%\begin{figure}[H]
%    \centering
%    \includegraphics[width=0.8\textwidth]{3-1.png}
%    \caption{Struktur file CRUD Tugas pada project Laravel}
%    \label{fig:3-1}
%\end{figure}

\section{Konfigurasi Database untuk Testing}

Pengujian dikonfigurasi menggunakan SQLite in-memory agar dapat berjalan di GitHub Actions tanpa memerlukan database server eksternal. Konfigurasi tersebut ditambahkan pada \texttt{phpunit.xml}.

\begin{lstlisting}[language=xml, caption={Konfigurasi SQLite in-memory di phpunit.xml}]
<env name="DB_CONNECTION" value="sqlite"/>
<env name="DB_DATABASE" value=":memory:"/>
\end{lstlisting}

Pengujian fitur dibuat pada \texttt{tests/Feature/TugasTest.php} mencakup tiga skenario: endpoint mengembalikan HTTP 200, pembuatan data baru, dan pengambilan data berdasarkan ID.

% Screenshot 2: Test file
%\begin{figure}[H]
%    \centering
%    \includegraphics[width=0.8\textwidth]{3-2.png}
%    \caption{File pengujian TugasTest.php}
%    \label{fig:3-2}
%\end{figure}

\section{Workflow CI/CD Empat Tahap}

Workflow CI/CD dikonfigurasi pada \texttt{.github/workflows/ci.yml} dengan empat job yang dirangkai menggunakan \texttt{needs:}.

\begin{lstlisting}[language=yaml, caption={Workflow CI/CD empat job}]
name: Laravel CI/CD Pipeline

on:
  push:
    branches: ['**']
  pull_request:
    branches: [main, dev]

jobs:
  build:
    runs-on: ubuntu-latest
    steps:
      - uses: actions/checkout@v4
      - uses: shivammathur/setup-php@v2
        with:
          php-version: '8.4'
          extensions: mbstring, dom, fileinfo, sqlite3
      - name: Cache Composer dependencies
        uses: actions/cache@v4
        with:
          path: vendor
          key: composer-${{ hashFiles('composer.lock') }}
      - name: Install dependencies
        run: composer install --no-interaction --prefer-dist
      - name: Upload vendor artifact
        uses: actions/upload-artifact@v4
        with:
          name: vendor
          path: vendor/

  test:
    runs-on: ubuntu-latest
    needs: build
    steps:
      - uses: actions/checkout@v4
      - uses: shivammathur/setup-php@v2
        with:
          php-version: '8.4'
          extensions: mbstring, dom, fileinfo, sqlite3
      - name: Download vendor artifact
        uses: actions/download-artifact@v4
        with:
          name: vendor
          path: vendor/
      - name: Prepare environment
        run: |
          cp .env.example .env
          php artisan key:generate
          touch database/database.sqlite
      - name: Run tests
        run: php artisan test

  staging:
    runs-on: ubuntu-latest
    needs: test
    steps:
      - uses: actions/checkout@v4
      - name: Deploy to staging
        run: |
          echo "=== STAGING DEPLOYMENT ==="
          echo "Branch: ${{ github.ref_name }}"
          echo "Commit: ${{ github.sha }}"
          echo "Simulating deploy to staging server..."
          echo "Staging deploy complete!"

  production:
    runs-on: ubuntu-latest
    needs: staging
    if: github.ref == 'refs/heads/main'
    environment: production
    steps:
      - uses: actions/checkout@v4
      - name: Execute deploy steps
        run: bash deploy.sh
\end{lstlisting}

Alur pipeline yang terbentuk adalah sebagai berikut:

\begin{lstlisting}
build --> test --> staging --> production (main only)
\end{lstlisting}

% Screenshot 3: Diagram pipeline
%\begin{figure}[H]
%    \centering
%    \includegraphics[width=0.8\textwidth]{3-3.png}
%    \caption{Tampilan pipeline empat tahap di GitHub Actions}
%    \label{fig:3-3}
%\end{figure}

\section{Proteksi Job Production}

Job \texttt{production} dilindungi dengan dua mekanisme:
\begin{itemize}
    \item Kondisi \texttt{if: github.ref == 'refs/heads/main'} memastikan job hanya berjalan ketika push dilakukan ke branch \texttt{main}.
    \item Penggunaan \texttt{environment: production} mengaktifkan mekanisme \textit{required reviewer} yang dapat dikonfigurasi pada pengaturan repository GitHub.
\end{itemize}

Konfigurasi environment \texttt{production} dilakukan melalui menu \textit{Settings} $\rightarrow$ \textit{Environments} pada repository GitHub.

% Screenshot 4: Environment production
%\begin{figure}[H]
%    \centering
%    \includegraphics[width=0.8\textwidth]{3-4.png}
%    \caption{Konfigurasi environment production dengan required reviewer}
%    \label{fig:3-4}
%\end{figure}

\section{Skrip Deploy (deploy.sh)}

Skrip \texttt{deploy.sh} disimpan di root repository dengan tujuh langkah deploy yang ditulis sebagai perintah \texttt{echo} untuk simulasi. Penggunaan \texttt{set -e} memastikan skrip berhenti jika salah satu langkah gagal.

\begin{lstlisting}[language=bash, caption={Isi deploy.sh dengan tujuh langkah}]
#!/bin/bash
set -e

echo "Step 1: Pull latest code from repository"
git pull origin main

echo "Step 2: Install/update Composer dependencies"
composer install --no-dev --optimize-autoloader

echo "Step 3: Run database migrations"
php artisan migrate --force

echo "Step 4: Clear and cache application config"
php artisan config:cache

echo "Step 5: Clear and cache routes"
php artisan route:cache

echo "Step 6: Clear and cache views"
php artisan view:cache

echo "Step 7: Restart queue workers and reload services"
php artisan queue:restart

echo "Deploy complete!"
\end{lstlisting}

% Screenshot 5: deploy.sh
%\begin{figure}[H]
%    \centering
%    \includegraphics[width=0.8\textwidth]{3-5.png}
%    \caption{File deploy.sh pada repository}
%    \label{fig:3-5}
%\end{figure}

\section{Pembuktian: Push ke Branch Fitur}

Untuk membuktikan bahwa job \texttt{production} dilewati pada branch selain \texttt{main}, dilakukan push ke branch \texttt{feature/crud-tugas}. Hasil pipeline menunjukkan job \texttt{build}, \texttt{test}, dan \texttt{staging} berhasil berjalan, sedangkan job \texttt{production} berstatus \textit{skipped}.

% Screenshot 6: Pipeline feature branch
%\begin{figure}[H]
%    \centering
%    \includegraphics[width=0.8\textwidth]{3-6.png}
%    \caption{Pipeline pada branch fitur: production skipped}
%    \label{fig:3-6}
%\end{figure}

\section{Pembuktian: Test Gagal Menghentikan Pipeline}

Untuk membuktikan pipeline berhenti ketika pengujian gagal, satu assertion pada test sengaja diubah menjadi salah. Setelah push, pipeline menunjukkan job \texttt{test} berstatus merah dan job selanjutnya (\texttt{staging} dan \texttt{production}) tidak dijalankan.

% Screenshot 7: Pipeline merah
%\begin{figure}[H]
%    \centering
%    \includegraphics[width=0.8\textwidth]{3-7.png}
%    \caption{Pipeline berhenti merah karena test gagal}
%    \label{fig:3-7}
%\end{figure}

Setelah test diperbaiki dan dipush kembali, pipeline kembali berjalan hijau seluruhnya.

% Screenshot 8: Pipeline hijau setelah perbaikan
%\begin{figure}[H]
%    \centering
%    \includegraphics[width=0.8\textwidth]{3-8.png}
%    \caption{Pipeline hijau setelah test diperbaiki}
%    \label{fig:3-8}
%\end{figure}

\section{Pipeline Production dari Branch Main}

Setelah Pull Request dari \texttt{feature/crud-tugas} ke \texttt{dev} dan dari \texttt{dev} ke \texttt{main} berhasil di-merge, pipeline dijalankan dari branch \texttt{main}. Seluruh empat job berhasil berjalan termasuk job \texttt{production}.

% Screenshot 9: Pipeline main hijau semua
%\begin{figure}[H]
%    \centering
%    \includegraphics[width=0.8\textwidth]{3-9.png}
%    \caption{Seluruh pipeline berhasil dari branch main termasuk production}
%    \label{fig:3-9}
%\end{figure}

%=== BAB 3 ===
\chapter{Kesimpulan}

Berdasarkan praktikum yang telah dilakukan, dapat disimpulkan beberapa hal berikut:
\begin{enumerate}
    \item Pipeline CI/CD empat tahap (\textit{build}, \textit{test}, \textit{staging}, \textit{production}) berhasil diimplementasikan menggunakan GitHub Actions dengan kata kunci \texttt{needs:} untuk mengatur dependensi antar-job.
    \item Job \texttt{build} berhasil memasang dependensi dan meneruskan artefak \texttt{vendor} ke job \texttt{test} menggunakan mekanisme \textit{upload/download artifact}.
    \item Proteksi job \texttt{production} berhasil diterapkan menggunakan kondisi \texttt{if: github.ref == 'refs/heads/main'} dan \texttt{environment: production}, sehingga job hanya berjalan dari branch \texttt{main}.
    \item Skrip \texttt{deploy.sh} dengan tujuh langkah deploy berhasil disimpan di repository dan dieksekusi oleh job \texttt{production}.
    \item Pipeline terbukti berhenti pada stage \texttt{test} ketika pengujian gagal, dan job \texttt{production} terbukti dilewati pada branch selain \texttt{main}.
\end{enumerate}

\clearpage

\begin{thebibliography}{9}
\bibitem{githubActionsDoc}
GitHub. (2026). \textit{GitHub Actions Documentation}. Diakses dari \url{https://docs.github.com/en/actions}.
\bibitem{laravelDoc}
Laravel. (2026). \textit{Laravel Documentation}. Diakses dari \url{https://laravel.com/docs}.
\end{thebibliography}

\end{document}
